% from ega3-1-fr.tex
\oldpage[III]{130}

un $\mathcal O_y$-module de type fini, donc séparé pour la topologie
$\mathfrak m_y$-préadique, et on a un isomorphisme topologique canonique
\[
\widehat{(R^pf_*(\mathcal F))_y}
\xrightarrow{\sim}
\varprojlim_kH^p\bigl(f^{-1}(y),
\mathcal F\otimes_{\mathcal O_Y}(\mathcal O_y/\mathfrak m_y^k)\bigr)
\tag{4.2.1.1}\label{III.4.2.1.1-fr}
\]
où le premier membre est le complété de $(R^pf_*(\mathcal F))_y$ pour la
topologie $\mathfrak m_y$-préadique, et au second membre $f^{-1}(y)$ est
considéré, pour tout $k\geq0$, comme espace sous-jacent au préschéma
$X\times_Y\operatorname{Spec}(\mathcal O_y/\mathfrak m_y^k)$
(I, 3.6.1).
\end{proposition}

\begin{proof}
Comme $\mathcal O_y$ est un anneau local noethérien et
$(R^pf_*(\mathcal F))_y$ un $\mathcal O_y$-module de type fini
(\hyperref[III.3.2.1-fr]{3.2.1}), la topologie $\mathfrak m_y$-préadique
sur $(R^pf_*(\mathcal F))_y$ est séparée (0\textsubscript{I}, 7.3.5). Les
autres assertions sont conséquences de (\hyperref[III.4.1.7-fr]{4.1.7})
lorsque $Y$ est noethérien et le point $y$ fermé, en remplaçant $Y$ par un
voisinage affine de $y$ et prenant $Y'=\{y\}$, compte tenu de
(G, II, 4.9.1). Dans le cas général, posons
\[
Y_1=\operatorname{Spec}(\mathcal O_y),\qquad
X_1=X\times_YY_1,\qquad
\mathcal F_1=\mathcal F\otimes_{\mathcal O_Y}\mathcal O_{Y_1}
\]
et soit $f_1=f\times1_{Y_1}:X_1\to Y_1$~; $Y_1$ est noethérien et $f_1$
propre (II, 5.4.2, (iii)) et $\mathcal F_1$ est cohérent
(0\textsubscript{I}, 5.3.11). Soit $y_1$ l'unique point fermé de $Y_1$~;
la proposition est valable pour $f_1$, $\mathcal F_1$ et $y_1$~; on a
$\mathcal O_{y_1}=\mathcal O_y$, $f_1^{-1}(y_1)=f^{-1}(y)$ (I, 3.6.5),
les préschémas
$X\times_Y\operatorname{Spec}(\mathcal O_y/\mathfrak m_y^k)$ et
$X_1\times_{Y_1}\operatorname{Spec}(\mathcal O_y/\mathfrak m_y^k)$ étant
canoniquement identifiés (I, 3.3.9)~; en outre,
$\mathcal F_1\otimes_{\mathcal O_{Y_1}}
(\mathcal O_y/\mathfrak m_y^k)$ s'identifie à
$\mathcal F\otimes_{\mathcal O_Y}(\mathcal O_y/\mathfrak m_y^k)$
(I, 9.1.6). Il reste à voir que $R^pf_{1*}(\mathcal F_1)$ est
canoniquement isomorphe à
$R^pf_*(\mathcal F)\otimes_{\mathcal O_Y}\mathcal O_{Y_1}$, ce qui
résulte de (\hyperref[III.1.4.15-fr]{1.4.15}), le morphisme local
$\operatorname{Spec}(\mathcal O_y)\to Y$ étant plat
(0\textsubscript{I}, 6.7.1 et I, 2.4.2).
\end{proof}

Le corollaire suivant utilise la terminologie de la théorie de la dimension
(chap. IV) et ne sera pas appliqué avant le chap. IV.

\begin{corollary}[4.2.2]
\label{III.4.2.2-fr}
Soient $Y$ un préschéma localement noethérien, $f:X\to Y$ un morphisme
propre, $y$ un point de $Y$, $r$ la dimension de $f^{-1}(y)$. Alors, pour
tout $\mathcal O_X$-Module cohérent $\mathcal F$, les faisceaux
$R^pf_*(\mathcal F)$ sont nuls au voisinage de $y$ pour tout $p>r$.
\end{corollary}

\begin{proof}
En effet, on a alors
\[
H^p\bigl(f^{-1}(y),
\mathcal F\otimes_{\mathcal O_Y}(\mathcal O_y/\mathfrak m_y^k)\bigr)=0
\]
(G, II, 4.15.2) pour tout $k$, donc
(\hyperref[III.4.2.1-fr]{4.2.1}) le séparé complété de
$(R^pf_*(\mathcal F))_y$ pour la topologie $\mathfrak m_y$-préadique est
nul, et comme cette topologie est séparée, on a aussi
$(R^pf_*(\mathcal F))_y=0$~; d'où la conclusion, puisque
$R^pf_*(\mathcal F)$ est cohérent (0\textsubscript{I}, 5.2.2).
\end{proof}

\begin{env}[4.2.3]
\label{III.4.2.3-fr}
Le résultat (\hyperref[III.4.2.1-fr]{4.2.1}) est surtout employé pour
$p=0$~; on obtient donc le corollaire suivant~:
\end{env}

\begin{corollary}[4.2.4]
\label{III.4.2.4-fr}
Sous les hypothèses de (\hyperref[III.4.2.1-fr]{4.2.1}), on a un
isomorphisme canonique topologique
\[
\widehat{(f_*(\mathcal F))_y}
\xrightarrow{\sim}
\varprojlim_k\Gamma\bigl(f^{-1}(y),
\mathcal F_y/\mathfrak m_y^k\mathcal F_y\bigr).
\]
\end{corollary}

\subsection{Le théorème de connexion de Zariski}
\label{subsection:III.4.3-fr}
\label{III.4.3-fr}

Les résultats de ce numéro et du suivant généralisent des théorèmes bien
connus de Zariski, et peuvent tous se déduire de
(\hyperref[III.4.2.4-fr]{4.2.4}). Ils sont conséquences du th. suivant~:

\begin{theorem}[4.3.1]
\label{III.4.3.1-fr}
\textnormal{(théorème de connexion).}
Soient $Y$ un préschéma localement noethérien, $f:X\to Y$ un morphisme
propre. Alors $\mathcal A(X)=f_*(\mathcal O_X)$ est une
$\mathcal O_Y$-Algèbre cohérente.

